\section{Scope and Summary}

\subsection{Question answered by this paper}

This paper explains what the selected \ProductReleaseTag{} implementation does, how its dual-branch
simulation is constructed, which statistical quantities its exact release evidence contains, how
the evidence is bound to code and runtime identities, and where the evidence stops. It is intended
for technical evaluators, model-risk teams, fair-outcomes specialists, security reviewers, and
procurement engineers who need more than a sample report.

The subject is one exact release: annotated tag \ProductReleaseTag{} at product commit
\IdentityText{\ProductCommitRaw}. The evidence source is release-evidence run \EvidenceRunRaw.
The intake and robustness producers completed successfully in retained attempt
\EvidenceSourceAttemptRaw; the overall workflow attempt later ended in failure. Current attempt
\EvidenceCurrentAttemptRaw{} succeeded. This split-attempt history is explicit because resolving
artifacts from the current attempt would select the wrong source.

\subsection{What FL-BSA is}

FL-BSA is a self-hosted fair-outcomes evidence appliance. It ingests a configured tabular credit
dataset inside the operator's environment, constructs two synthetic simulation branches, evaluates
disparity and quality diagnostics, and emits reports, manifests, certificates, and evidence bundles.
It does not approve, deny, rank, or override live credit decisions. The default v5 workflow is a
simulation path rather than a production scoring service.

The selected v5 backend is \texttt{first\_party\_evidence\_native}. It is a profile-based tabular
generator, not CTGAN, DataCebo, a GAN, or a neural generator. CTGAN remains an explicitly gated
legacy comparison path and is outside the evidence characterized here.

\subsection{Release-bound findings}

\begin{KeyBox}
\begin{KeyValueList}
  \item[Gender disparity] Amplification AIR is \GenderAIRAmplification{}
    [\GenderAIRAmplificationLCI{}, \GenderAIRAmplificationUCI{}] for female versus male selection.
    The point estimate and interval are below the internal \AIRThreshold{} screen.
  \item[Branch separation] Historical-fixture AIR is \GenderAIRHistorical{}; amplification AIR is
    \GenderAIRAmplification{}; the intrinsic equal-rate policy control is \GenderAIRIntrinsic{}.
  \item[Availability] AOD and EOD are unavailable. Equalized-odds results are not informative on the
    degenerate confusion surface. ECE is not evaluated because the scenario has no probability-score
    surface.
  \item[Quality] Composite quality is \AmpQuality{} for amplification and \IntQuality{} for
    intrinsic against a configured 0.80 threshold. These are quality composites, not utility or
    privacy guarantees.
  \item[Robustness] \RobustnessPasses{}/\RobustnessRuns{} scenario contracts passed across four
    scenarios and ten predeclared seeds. A scenario-contract pass can include an expected fairness
    warning, as in the gender-bias positive control.
  \item[Utility] Amplification advisory AUC ranges from \UtilityMinAUC{} to \UtilityMaxAUC{}, with
    mean \UtilityMeanAUC{}. \UtilityBelowFloor{} of 40 runs fall below the 0.70 floor. Accuracy
    ranges from \UtilityMinAccuracy{} to \UtilityMaxAccuracy{}, with mean \UtilityMeanAccuracy{};
    it has no configured threshold. Intrinsic utility was not evaluated. Production utility is not
    established.
\end{KeyValueList}
\end{KeyBox}

\subsection{Claim boundary}

This paper and companion have disposition \texttt{characterization\_only} and are
customer-evidence ineligible. The underlying product release evidence has separately scoped native
eligibility metadata; that eligibility does not transfer to these public artifacts. The reviewed
correction changes only the whitepaper field. Public availability does not convert this
synthetic/demo characterization into customer validation, production qualification, certification,
compliance, regulator approval, Marketplace/GA availability, or a claim that the method improves
real lending outcomes.
